\documentclass[10pt]{article}
\usepackage[T1]{fontenc}
\usepackage[utf8]{inputenc}
\usepackage[french]{babel}
\usepackage{amsmath,amssymb,mathrsfs,mathpazo}
\usepackage{geometry}
\geometry{a4paper,margin=2cm}
\usepackage{graphicx,array,tabularx,tabularray,colortbl,multirow,multicol,booktabs,enumitem,pifont,ulem,fancybox,diagbox,cancel}
\usepackage{pstricks,pstricks-add,pst-plot,pst-tree,pst-node,pst-text,pst-func,pst-math,pst-bspline,pst-3dplot,pst-eucl,pst-solides3d}
\let\pscalculate\undefined
\usepackage{pst-bezier}
\usepackage{tikz,pgfplots,tkz-tab}
\usetikzlibrary{arrows,arrows.meta,patterns,calc,positioning,trees,decorations.markings,shapes.geometric,shapes.misc,shapes.arrows,intersections,decorations.pathmorphing,decorations.pathreplacing,backgrounds,shadows,babel}
\pgfplotsset{compat=1.18}
\usepgfplotslibrary{fillbetween}
\usepackage[np]{numprint}
\usepackage{listings,makecell,pas-tableur,twoopt}
\let\barre\undefined
\usepackage[pstricks]{tablvar}
\definecolor{darkgreen}{RGB}{0,100,0}
\providecommand{\red}{\color{red}}
\providecommand{\blue}{\color{blue}}
\providecommand{\green}{\color{green}}
\providecommand{\black}{\color{black}}
\providecommand{\white}{\color{white}}
\providecommand{\gray}{\color{gray}}
\providecommand{\cyan}{\color{cyan}}
\providecommand{\magenta}{\color{magenta}}
\providecommand{\yellow}{\color{yellow}}
\providecommand{\violet}{\color{violet}}
\providecommand{\orange}{\color{orange}}
\providecommand{\euro}{\texteuro}
\providecommand{\EUR}{\texteuro}
\providecommand{\ieme}{\textsuperscript{e}}
\providecommand{\degres}{\ensuremath{^\circ}}
\providecommand{\degree}{\ensuremath{^\circ}}
\providecommand{\vv}[1]{\overrightarrow{#1}}
\providecommand{\Coord}[1]{#1}
\providecommand{\decosix}{}
\providecommand{\starredbullet}{\ensuremath{\bullet}}
\pagestyle{empty}
\setlength{\parindent}{0pt}
\usepackage{adjustbox}\DeclareUnicodeCharacter{2236}{:}\DeclareUnicodeCharacter{2190}{\ensuremath{\leftarrow}}
\providecommand{\R}{\mathbb{R}}
\providecommand{\N}{\mathbb{N}}
\providecommand{\Z}{\mathbb{Z}}
\providecommand{\Q}{\mathbb{Q}}
\providecommand{\D}{\mathbb{D}}
\providecommand{\C}{\mathbb{C}}
\providecommand{\e}{{\mathrm{e}}}
\providecommand{\vect}[1]{\overrightarrow{#1}}
\providecommand{\vectt}[1]{\overrightarrow{\mathrm{#1}}}
\providecommand{\vvt}[1]{\overrightarrow{\mathrm{#1}}}
\providecommand{\np}[2][]{{#2~#1}}
\providecommand{\vv}[1]{{\overrightarrow{#1}}}
\renewcommand{\d}{{\mathrm{d}}}
\providecommand{\degres}{{^\circ}}
\providecommand{\degre}{{^\circ}}
\providecommand{\blue}{\color{blue}}
\providecommand{\red}{\color{red}}
\providecommand{\enskip}{\quad}
\providecommand{\llbracket}{\lbrack\!\lbrack}
\providecommand{\rrbracket}{\rbrack\!\rbrack}
\providecommand{\euro}{{\text{€}}}
\providecommand{\no}{{\text{n°}}}
\providecommand{\interval}[2]{{[#1;#2]}}
\providecommand{\intervFO}[2]{{[#1;#2[}}
\providecommand{\dotfill}{\ldots}

\providecommand{\ds}{\displaystyle}
\providecommand{\textding}[1]{#1}
\providecommand{\Oij}{\(\mathrm{O};\vec{\imath},\vec{\jmath}\)}
\providecommand{\Oijk}{\(\mathrm{O};\vec{\imath},\vec{\jmath},\vec{k}\)}
\providecommand{\Ouv}{\(\mathrm{O};\vec{u},\vec{v}\)}
\providecommand{\intervFF}[2]{\ensuremath{[#1;#2]}}
\renewcommand{\euro}{\ifmmode\text{\texteuro}\else\texteuro\fi}
\DeclareUnicodeCharacter{00B0}{\ifmmode^{\circ}\else\textdegree\fi}
\DeclareUnicodeCharacter{20AC}{\ifmmode\text{\texteuro}\else\texteuro\fi}
% Énoncés : Éducation nationale. Transcription et corrigés : APMEP, Denis Vergès et contributeurs. Découpage et classement : CoopMaths, Sébastien Lozano, Éric Elter et contributeurs. Adaptation de la présentation : Maths Explorer.
% CC BY-SA : https://coopmaths.fr/www/about/
% https://coopmaths.fr/alea/static/bac/2021/tex/bac_2021_03_sujet2_2.tex
\begin{document}
\noindent\textit{APMEP / CoopMaths. CC BY-SA. Présentation adaptée par Maths Explorer.}\par\medskip
\medskip

On considère les suites $\left(u_n\right)$ et $\left(v_n\right)$ définies pour tout entier naturel $n$ par : 

\[\left\{\begin{array}{l !{=} l}
u_0		& v_0 = 1\\
u_{n+1}	&  u_n +v_n\\
v_{n+1}	&  2u_n + v_n
\end{array}\right.\]

Dans toute la suite de l'exercice, on \textbf{admet} que les suites $\left(u_n\right)$ et $\left(v_n\right)$ \textbf{sont strictement positives}.

\medskip

\begin{enumerate}
\item 
	\begin{enumerate}
		\item Calculez $u_1$ et $v_1$.
		\item Démontrer que la suite $\left(v_n\right)$ est strictement croissante, puis en déduire que, pour tout entier naturel $n$,\: $v_n \geqslant 1$.
		\item Démontrer par récurrence que, pour tout entier naturel $n$, on a : $u_n \geqslant  n + 1$.
		\item En déduire la limite de la suite $\left(u_n\right)$.
	\end{enumerate}
\item On pose, pour tout entier naturel $n$ :

\[r_n = \dfrac{v_n}{u_n}.\]

On admet que:

\[r_n^2 = 2 + \dfrac{(- 1)^{n+1}}{u_n^2}\]

	\begin{enumerate}
		\item Démontrer que pour tout entier naturel $n$ :

\[- \dfrac{1}{u_n^2} \leqslant \dfrac{(- 1)^{n+1}}{u_n^2} \leqslant \dfrac{1}{u_n^2}.\]
		\item En déduire :

\[\displaystyle\lim_{n \to + \infty} \dfrac{(- 1)^{n+1}}{u_n^2}.\]

		\item Déterminer la limite de la suite $\left(r_n^2\right)$ et en déduire que $\left(r_n\right)$ converge vers $\sqrt{2}$.
		\item  Démontrer que pour tout entier naturel $n$,

\[r_{n+1} = \dfrac{2 + r_n}{1 + r_n}.\]

		\item On considère le programme suivant écrit en langage Python :
		
\begin{center}
\fbox{
\begin{tabularx}{0.5\linewidth}{X}
\textbf{def seuil()}:\\
\qquad  n = 0\\
\qquad  r = 1\\
\qquad \textbf{while} abs(r-sqrt(2)) > 10**(-4) :\\
\quad \qquad r = (2+r)/(1+r)\\
\quad \qquad n = n+1\\
\qquad \textbf{return} n\\ 
\end{tabularx}
}
\end{center}

\smallskip

(abs désigne la valeur absolue, sqrt la racine carrée et 10** (-4) représente $10^{-4}$).

La valeur de $n$ renvoyée par ce programme est 5.

 À quoi correspond-elle ?
	\end{enumerate}
\end{enumerate}

\bigskip

\end{document}
